\documentclass[11pt]{article}

\usepackage[margin=2.5cm]{geometry}
\usepackage{amsmath,amssymb,amsthm}
\usepackage{bm}
\usepackage{enumitem}

\title{Correction -- Exercise Sheet\\[0.3em]
{\large Probability Spaces \& Random Variables}}
\author{}
\date{}

\begin{document}
\maketitle

\noindent
This document provides detailed solutions, exercise by exercise, for the sheet.

\bigskip
\hrule
\bigskip

% ============================================================
\section*{Part A -- Fundamental Exercises}
\begin{enumerate}

% -----------------------------
\item \textbf{Sets, unions/intersections, and $\limsup$ / $\liminf$}

\begin{enumerate}[label=(\alph*)]

\item  \emph{Dice roll}

\textbf{Correction}

Let $\Omega=\{1,2,3,4,5,6\}$,
\[
A=\{2\},\qquad B=\{2,4,6\},\qquad C=\{1,2,3\}.
\]

\medskip
\textbf{Complements.}
\[
A^c=\Omega\setminus A=\{1,3,4,5,6\},
\]
\[
B^c=\Omega\setminus B=\{1,3,5\},
\]
\[
C^c=\Omega\setminus C=\{4,5,6\}.
\]

\medskip
\textbf{Union, intersection, and difference.}
\[
A\cup C=\{2\}\cup\{1,2,3\}=\{1,2,3\},
\]
\[
B\cap C=\{2,4,6\}\cap\{1,2,3\}=\{2\},
\]
\[
B\setminus C=\{2,4,6\}\setminus\{1,2,3\}=\{4,6\}.
\]

\medskip
\textbf{Verification of $B\setminus C=B\cap C^c$.}
Since $C^c=\{4,5,6\}$, we have
\[
B\cap C^c=\{2,4,6\}\cap\{4,5,6\}=\{4,6\}.
\]
Thus,
\[
B\setminus C=\{4,6\}=B\cap C^c.
\]

\item \emph{Compute $\bigcup_{n\ge1}A_n$ and $\bigcap_{n\ge1}A_n$.}
For $n\ge1$,
\[
A_n=\Big[-\frac1n,\;2+\frac1n\Big],
\qquad
B_n=\Big[-\frac5n,\;\frac{n}{2}\Big].
\]


Note that $(A_n)$ is \textbf{decreasing} (nested): indeed
\[
-\frac{1}{n+1} > -\frac{1}{n}
\quad\text{and}\quad
2+\frac{1}{n+1} < 2+\frac{1}{n}
\;\;\Rightarrow\;\;
A_{n+1}\subset A_n.
\]
Hence
\[
\bigcup_{n\ge1} A_n = A_1 = [-1,3],
\qquad
\bigcap_{n\ge1} A_n
=
\left[\lim_{n\to\infty}\left(-\frac1n\right),\,\lim_{n\to\infty}\left(2+\frac1n\right)\right]
=
[0,2].
\]

\item \emph{Compute $\limsup A_n$ and $\liminf A_n$ and interpret them.}

Recall
\[
\liminf_{n\to\infty}A_n=\bigcup_{n\ge1}\bigcap_{k\ge n}A_k,
\qquad
\limsup_{n\to\infty}A_n=\bigcap_{n\ge1}\bigcup_{k\ge n}A_k.
\]
For a \textbf{monotone decreasing} sequence, we have
\[
\liminf A_n=\limsup A_n=\bigcap_{n\ge1}A_n.
\]
Therefore,
\[
\liminf_{n\to\infty}A_n=\limsup_{n\to\infty}A_n=[0,2].
\]
Interpretation:
\begin{itemize}
\item $\liminf A_n$ = points that are eventually in all $A_n$ (``eventually always'').
\item $\limsup A_n$ = points that belong to infinitely many $A_n$ (``infinitely often'').
\end{itemize}
Here both coincide because the sets shrink nicely to a limit interval.

\item \emph{Same questions for $(B_n)_{n\ge1}$.}

First compute union/intersection directly.

\underline{Union.}
If $x\ge -5$, then choose $n$ large enough so that $\frac{n}{2}\ge x$ (always possible since $n/2\to\infty$).
Also we need $x\ge -\frac{5}{n}$, which holds for some $n$:
\begin{itemize}
\item if $x\ge 0$, it holds for every $n$,
\item if $-5\le x<0$, it holds for all integers $n\le -5/x$ (so at least one exists).
\end{itemize}
Thus every $x\ge -5$ belongs to at least one $B_n$, and no $x<-5$ can belong to any $B_n$.
So
\[
\bigcup_{n\ge1} B_n = [-5,\infty).
\]

\underline{Intersection.}
We need $x\in B_n$ for all $n$, i.e.
\[
x\ge -\frac{5}{n}\ \forall n
\quad\text{and}\quad
x\le \frac{n}{2}\ \forall n.
\]
The first condition forces $x\ge \sup_n(-5/n)=0$. The second forces $x\le \inf_n(n/2)=1/2$.
Hence
\[
\bigcap_{n\ge1} B_n = [0,1/2].
\]

\underline{$\liminf$ and $\limsup$.}
Fix $x\in\mathbb{R}$.
\begin{itemize}
\item If $x<0$, then for $n$ large enough, $-\frac{5}{n}>x$, hence $x\notin B_n$ eventually. So $x$ occurs only finitely many times and $x\notin \limsup B_n$.
\item If $x\ge 0$, then $x\ge -5/n$ for all $n$, and also $x\le n/2$ for all sufficiently large $n$.
Thus $x\in B_n$ for all large $n$, i.e. $x\in \liminf B_n$.
\end{itemize}
Therefore,
\[
\liminf_{n\to\infty}B_n=\limsup_{n\to\infty}B_n=[0,\infty).
\]
\end{enumerate}

\bigskip

% -----------------------------
\item \textbf{Basic probability measure properties (quick proofs)}

Let $(\Omega,\mathcal{F},\mathbb{P})$ be a probability space and $A,B\in\mathcal{F}$.

\begin{enumerate}[label=(\alph*)]
\item If $A\subseteq B$, then $B=A\cup(B\setminus A)$ with $A$ and $(B\setminus A)$ disjoint.
So by additivity,
\[
\mathbb{P}(B)=\mathbb{P}(A)+\mathbb{P}(B\setminus A)\ge \mathbb{P}(A).
\]

\item Since $\Omega=A\cup A^c$ and $A\cap A^c=\emptyset$,
\[
1=\mathbb{P}(\Omega)=\mathbb{P}(A)+\mathbb{P}(A^c)\quad\Rightarrow\quad \mathbb{P}(A^c)=1-\mathbb{P}(A).
\]

\item Decompose $A\cup B$ into disjoint pieces:
\[
A\cup B = A \cup (B\setminus A),
\quad\text{and}\quad
\mathbb{P}(B)=\mathbb{P}(B\setminus A)+\mathbb{P}(A\cap B).
\]
Hence
\[
\mathbb{P}(A\cup B)=\mathbb{P}(A)+\mathbb{P}(B\setminus A)
= \mathbb{P}(A)+\mathbb{P}(B)-\mathbb{P}(A\cap B).
\]

\item Since $\mathbb{P}(A\cap B)\ge 0$,
\[
\mathbb{P}(A\cup B)=\mathbb{P}(A)+\mathbb{P}(B)-\mathbb{P}(A\cap B)\le \mathbb{P}(A)+\mathbb{P}(B).
\]


\medskip

% -----------------------------
\item \textbf{Union bound practice }

\emph{Claim.} For events $A_1,\dots,A_n$,
\[
\mathbb{P}\Big(\bigcup_{k=1}^n A_k\Big)\le \sum_{k=1}^n \mathbb{P}(A_k).
\]

\emph{Proof by induction.}
For $n=2$, it is the union bound from Part (A2d).
Assume it holds for $n-1$. Then
\[
\mathbb{P}\Big(\bigcup_{k=1}^n A_k\Big)
=
\mathbb{P}\Big(\Big(\bigcup_{k=1}^{n-1}A_k\Big)\cup A_n\Big)
\le
\mathbb{P}\Big(\bigcup_{k=1}^{n-1}A_k\Big)+\mathbb{P}(A_n)
\le
\sum_{k=1}^{n}\mathbb{P}(A_k).
\]
\bigskip
\end{enumerate}
% -----------------------------
\item \textbf{Pairwise independence vs mutual independence}

$\Omega=\{\omega_1,\omega_2,\omega_3,\omega_4\}$ with uniform probability, so each atom has prob.\ $1/4$.
\[
A=\{\omega_1,\omega_2\},\quad
B=\{\omega_1,\omega_3\},\quad
C=\{\omega_2,\omega_3\}.
\]

\begin{enumerate}[label=(\alph*)]
\item $\mathbb{P}(A)=\mathbb{P}(B)=\mathbb{P}(C)=2/4=1/2$.
Also $A\cap B=\{\omega_1\}$ so $\mathbb{P}(A\cap B)=1/4=(1/2)(1/2)$.
Similarly $A\cap C=\{\omega_2\}$ and $B\cap C=\{\omega_3\}$, each with probability $1/4$.
Hence $A,B,C$ are pairwise independent.

\item $A\cap B\cap C = \emptyset$ (no outcome is in all three sets), so
\[
\mathbb{P}(A\cap B\cap C)=0,
\qquad
\mathbb{P}(A)\mathbb{P}(B)\mathbb{P}(C)=(1/2)^3=1/8.
\]
Not equal, so they are not mutually independent.

\item Lesson: \textbf{pairwise independence does not imply mutual independence.}
\end{enumerate}

\bigskip

% -----------------------------
\item \textbf{Independence (coin tosses)}

Three independent fair tosses $(X_1,X_2,X_3)\in\{H,T\}^3$.
\[
A=\{X_1=H\},\qquad
B=\{X_2=H\},\qquad
C=\{\text{exactly two heads occur}\}.
\]

\begin{enumerate}[label=(\alph*)]
\item $\mathbb{P}(A)=\mathbb{P}(B)=1/2$ and by independence of tosses
\[
\mathbb{P}(A\cap B)=\mathbb{P}(X_1=H,X_2=H)=1/4=\mathbb{P}(A)\mathbb{P}(B).
\]

\item Compute $\mathbb{P}(C)=3/8$ (the outcomes are HHT, HTH, THH).
Also $A\cap C$ means $X_1=H$ and exactly two heads total, so $(X_2,X_3)$ must contain exactly one head:
\[
A\cap C=\{\text{HHT or HTH}\}\quad\Rightarrow\quad \mathbb{P}(A\cap C)=2/8=1/4.
\]
But $\mathbb{P}(A)\mathbb{P}(C)=(1/2)(3/8)=3/16\neq 1/4$, hence $A$ and $C$ are \textbf{not} independent.

\item Example of dependence: let $D=\{X_1=X_2\}$. Then $\mathbb{P}(D)=1/2$ but
\[
\mathbb{P}(A\cap D)=\mathbb{P}(X_1=H,X_2=H)=1/4\neq (1/2)(1/2)=1/4
\]
(here it happens to match; choose instead $D'=\{X_1=X_2=X_3\}$: then $\mathbb{P}(D')=1/4$ and
$\mathbb{P}(A\cap D')=\mathbb{P}(\text{HHH})=1/8\neq (1/2)(1/4)=1/8$ again matches).
A clearer dependent pair is $A$ and $E=\{\text{at least one head}\}$:
$\mathbb{P}(E)=7/8$, but $\mathbb{P}(A\cap E)=\mathbb{P}(A)=1/2\neq (1/2)(7/8)$.
\end{enumerate}

\bigskip

% -----------------------------
\item \textbf{Conditioning: the two-children problem}

Equiprobable configurations:
$(F,F),(F,G),(G,F),(G,G)$.

\begin{enumerate}[label=(\alph*)]
\item Condition ``younger is a girl'' means second coordinate $F$:
possible outcomes $(F,F)$ and $(G,F)$, 2 equally likely cases, 1 favorable.
So
\[
\mathbb{P}(\text{both girls}\mid \text{younger is a girl})=\frac{1}{2}.
\]

\item Condition ``older is a girl'' means first coordinate $F$:
possible outcomes $(F,F)$ and $(F,G)$, hence again
\[
\mathbb{P}(\text{both girls}\mid \text{older is a girl})=\frac{1}{2}.
\]

\item Condition ``at least one is a girl'' excludes only $(G,G)$, leaving 3 equally likely cases.
Only $(F,F)$ is favorable, so
\[
\mathbb{P}(\text{both girls}\mid \text{at least one is a girl})=\frac{1}{3}.
\]
\end{enumerate}

\bigskip

% -----------------------------
\item \textbf{Law of total probability + Bayes (finite case)}

$Y\in\{G,B\}$ with $\mathbb{P}(Y=G)=0.7$, $\mathbb{P}(Y=B)=0.3$.
Signal $S\in\{1,0\}$ with
\[
\mathbb{P}(S=1\mid Y=G)=0.9,
\qquad
\mathbb{P}(S=1\mid Y=B)=0.2.
\]

\begin{enumerate}[label=(\alph*)]
\item Total probability:
\[
\mathbb{P}(S=1)=\mathbb{P}(S=1\mid Y=G)\mathbb{P}(Y=G)+\mathbb{P}(S=1\mid Y=B)\mathbb{P}(Y=B)
=0.9\cdot0.7+0.2\cdot0.3=0.69.
\]

\item Bayes:
\[
\mathbb{P}(Y=G\mid S=1)=\frac{\mathbb{P}(S=1\mid Y=G)\mathbb{P}(Y=G)}{\mathbb{P}(S=1)}
=
\frac{0.9\cdot0.7}{0.69}
=
\frac{0.63}{0.69}
=
\frac{21}{23}\approx 0.913.
\]

\item Interpretation: prior belief $0.7$ increases to about $0.913$ after seeing $S=1$.
\end{enumerate}

\bigskip

% -----------------------------
\item \textbf{Bernoulli/binomial modeling (oil slicks)}

Probability of oil slick per drilling: $p=0.1$.

\begin{enumerate}[label=(\alph*)]
\item One drilling $\leadsto$ Bernoulli r.v. $I\in\{0,1\}$ where $I=1$ if slick occurs.
Then $\mathbb{P}(I=1)=p$.

\item For $10$ drillings, $X=\#$ of slicks satisfies $X\sim\mathrm{Bin}(10,0.1)$ if we assume:
\begin{itemize}
\item identical probability $p=0.1$ each time,
\item independence between drillings.
\end{itemize}

\item If $X\sim\mathrm{Bin}(10,0.1)$:
\[
\mathbb{P}(X=2)=\binom{10}{2}(0.1)^2(0.9)^8,
\qquad
\mathbb{P}(X\ge 1)=1-\mathbb{P}(X=0)=1-(0.9)^{10}.
\]
\end{enumerate}

\bigskip

% -----------------------------
\item \textbf{CDF from a discrete distribution}

\[
\mathbb{P}(X=-1)=\frac14,\quad
\mathbb{P}(X=0)=\frac12,\quad
\mathbb{P}(X=2)=\frac14.
\]

\begin{enumerate}[label=(\alph*)]
\item The cdf is piecewise constant with jumps at $-1,0,2$:
\[
F_X(t)=\mathbb{P}(X\le t)=
\begin{cases}
0, & t<-1,\\[0.2em]
\frac14, & -1\le t<0,\\[0.2em]
\frac34, & 0\le t<2,\\[0.2em]
1, & t\ge 2.
\end{cases}
\]

\item It is non-decreasing and right-continuous by inspection. Jumps:
\[
\Delta F(-1)=\frac14,\quad \Delta F(0)=\frac12,\quad \Delta F(2)=\frac14.
\]

\item Expectation:
\[
\mathbb{E}[X]=(-1)\frac14+0\cdot\frac12+2\cdot\frac14=\frac14.
\]
Second moment:
\[
\mathbb{E}[X^2]=1\cdot\frac14+0+4\cdot\frac14=\frac54.
\]
Variance:
\[
\mathrm{Var}(X)=\mathbb{E}[X^2]-\mathbb{E}[X]^2=\frac54-\left(\frac14\right)^2
=\frac{20}{16}-\frac{1}{16}=\frac{19}{16}.
\]
\end{enumerate}

\bigskip

% -----------------------------
\item \textbf{Mean/variance for standard laws (short answers)}

\begin{enumerate}[label=(\alph*)]
\item If $X\sim\mathrm{Bin}(n,p)$:
\[
\mathbb{E}[X]=np,\qquad \mathrm{Var}(X)=np(1-p).
\]

\item If $Y\sim\mathrm{Poi}(\lambda)$:
\[
\mathbb{E}[Y]=\lambda,\qquad \mathrm{Var}(Y)=\lambda.
\]

\item If $U\sim U[a,b]$:
\[
\mathbb{E}[U]=\frac{a+b}{2},\qquad \mathrm{Var}(U)=\frac{(b-a)^2}{12}.
\]

\item If $E\sim \mathrm{Exp}(\lambda)$:
\[
\mathbb{E}[E]=\frac{1}{\lambda},\qquad \mathrm{Var}(E)=\frac{1}{\lambda^2}.
\]

\item If $Z\sim \mathcal{N}(\mu,\sigma^2)$:
\[
\mathbb{E}[Z]=\mu,\qquad \mathrm{Var}(Z)=\sigma^2.
\]
\end{enumerate}

\bigskip

% -----------------------------
\item \textbf{Transformation of a continuous r.v.: $X^2$}

Let $X\sim U[0,1]$ and $Z=X^2$.

\begin{enumerate}[label=(\alph*)]
\item For $t<0$, $F_Z(t)=0$. For $0\le t\le 1$,
\[
F_Z(t)=\mathbb{P}(X^2\le t)=\mathbb{P}(X\le \sqrt{t})=\sqrt{t}.
\]
For $t\ge 1$, $F_Z(t)=1$.

\item For $0<t<1$,
\[
f_Z(t)=F_Z'(t)=\frac{1}{2\sqrt{t}}\mathbf{1}_{(0,1)}(t).
\]

\item Directly: $\mathbb{E}[Z]=\mathbb{E}[X^2]=\int_0^1 x^2\,dx=1/3$.

Using $f_Z$:
\[
\mathbb{E}[Z]=\int_0^1 t\,\frac{1}{2\sqrt{t}}\,dt
=\frac12\int_0^1 \sqrt{t}\,dt
=\frac12\cdot\frac{2}{3}=\frac13.
\]
\end{enumerate}

\bigskip

% -----------------------------
\item \textbf{Censored exponential: mixed distribution}

Let $X\sim \mathrm{Exp}(1)$ and fix $a>0$. Define $Y=\min(X,a)$.

\begin{enumerate}[label=(\alph*)]
\item For $t<0$, $F_Y(t)=0$.
For $0\le t<a$,
\[
F_Y(t)=\mathbb{P}(\min(X,a)\le t)=\mathbb{P}(X\le t)=1-e^{-t}.
\]
For $t\ge a$, $\min(X,a)\le t$ always, so $F_Y(t)=1$.
Thus
\[
F_Y(t)=
\begin{cases}
0, & t<0,\\
1-e^{-t}, & 0\le t<a,\\
1, & t\ge a.
\end{cases}
\]
There is a jump at $t=a$ of size $\mathbb{P}(Y=a)=\mathbb{P}(X>a)=e^{-a}$.

\item $Y$ has a density on $[0,a)$ equal to $e^{-t}$, but also an \textbf{atom} at $a$:
\[
\mathbb{P}(Y=a)=e^{-a}>0.
\]
So $Y$ is \emph{mixed} (not purely continuous).

\item Using $Y=X\mathbf{1}_{\{X\le a\}}+a\mathbf{1}_{\{X>a\}}$:
\[
\mathbb{E}[Y]=\mathbb{E}[X\mathbf{1}_{X\le a}]+a\,\mathbb{P}(X>a).
\]
Compute
\[
\mathbb{E}[X\mathbf{1}_{X\le a}]
=\int_0^a t e^{-t}\,dt
=1-(a+1)e^{-a},
\qquad
\mathbb{P}(X>a)=e^{-a}.
\]
Hence
\[
\mathbb{E}[Y]=\big(1-(a+1)e^{-a}\big)+a e^{-a}=1-e^{-a}.
\]
\end{enumerate}

\bigskip

% -----------------------------
\item \textbf{Markov and Chebyshev: basic bounds}

\begin{enumerate}[label=(\alph*)]
\item If $X\ge0$ a.s.\ and $\mathbb{E}[X]=5$, Markov gives
\[
\mathbb{P}(X\ge 20)\le \frac{\mathbb{E}[X]}{20}=\frac{5}{20}=\frac14.
\]

\item If $\mathbb{E}[X]=2$ and $\mathrm{Var}(X)=9$, Chebyshev gives
\[
\mathbb{P}(|X-2|\ge 6)\le \frac{\mathrm{Var}(X)}{6^2}=\frac{9}{36}=\frac14.
\]

\item These bounds can be loose because they use only low-order moments and ignore the true tail shape.
\end{enumerate}

\end{enumerate}

\bigskip
\hrule
\bigskip

% ============================================================
\section*{Part B -- Advanced and Finance-Oriented Exercises}
\begin{enumerate}


% =========================================================
% Solutions (in LaTeX) — Probability Foundations & RV
% =========================================================

% -----------------------------
\item \textbf{Limsup/liminf of events and interpretation — Solution}

Recall
\[
\limsup_{n\to\infty}A_n=\bigcap_{n\ge1}\ \bigcup_{k\ge n} A_k,
\qquad
\liminf_{n\to\infty}A_n=\bigcup_{n\ge1}\ \bigcap_{k\ge n} A_k.
\]

\begin{enumerate}[label=(\alph*)]
\item \textbf{Show that $\liminf A_n \subseteq \limsup A_n$.}

Let $\omega\in\liminf A_n$. Then by definition there exists $n_0\ge1$ such that
\[
\omega\in \bigcap_{k\ge n_0}A_k,
\]
i.e.\ $\omega\in A_k$ for all $k\ge n_0$. In particular, for every $n\ge1$ we have
\[
\omega\in \bigcup_{k\ge n}A_k
\quad\text{(take any }k\ge \max\{n,n_0\}\text{)}.
\]
Hence $\omega\in\bigcap_{n\ge1}\bigcup_{k\ge n}A_k=\limsup A_n$, proving
$\liminf A_n\subseteq \limsup A_n$.

\item \textbf{Interpretation.}

\[
\limsup_{n\to\infty}A_n
=
\{\omega:\ \omega\in A_n\ \text{for infinitely many }n\}
\quad\text{(``infinitely often'')},
\]
\[
\liminf_{n\to\infty}A_n
=
\{\omega:\ \exists n_0,\ \forall n\ge n_0,\ \omega\in A_n\}
\quad\text{(``eventually always'')}.
\]

\item \textbf{Example with $\liminf A_n=\emptyset$ but $\limsup A_n\neq\emptyset$.}

Take $\Omega=\mathbb{N}$ and define
\[
A_n=
\begin{cases}
\{\omega_0\} & \text{if } n \text{ is even},\\
\emptyset & \text{if } n \text{ is odd},
\end{cases}
\]
for some fixed $\omega_0\in\Omega$.
Then $\omega_0\in A_n$ for infinitely many $n$ (all even $n$), so $\omega_0\in\limsup A_n$
and $\limsup A_n=\{\omega_0\}\neq\emptyset$.
But $\omega_0$ is \emph{not} eventually always in $A_n$ (odd indices fail), hence
$\liminf A_n=\emptyset$.
\end{enumerate}

% -----------------------------
\item \textbf{Monotone sequences of events (continuity of probability) — Solution}

Let $(A_n)$ be increasing and $(B_n)$ decreasing.

\begin{enumerate}[label=(\alph*)]
\item \textbf{Show $\mathbb{P}(\cup_{n\ge1}A_n)=\lim_{n\to\infty}\mathbb{P}(A_n)$.}

Define $A=\bigcup_{n\ge1}A_n$ and $C_1=A_1$, $C_n=A_n\setminus A_{n-1}$ for $n\ge2$.
Because $(A_n)$ is increasing, $(C_n)$ are pairwise disjoint and
\[
A=\bigcup_{n\ge1}C_n,
\qquad
A_m=\bigcup_{n=1}^m C_n.
\]
By countable additivity,
\[
\mathbb{P}(A)=\sum_{n\ge1}\mathbb{P}(C_n),
\qquad
\mathbb{P}(A_m)=\sum_{n=1}^m\mathbb{P}(C_n).
\]
Taking $m\to\infty$ yields
\[
\lim_{m\to\infty}\mathbb{P}(A_m)=\sum_{n\ge1}\mathbb{P}(C_n)=\mathbb{P}(A).
\]

\item \textbf{Show $\mathbb{P}(\cap_{n\ge1}B_n)=\lim_{n\to\infty}\mathbb{P}(B_n)$.}

Set $C_n=B_n^c$. Since $(B_n)$ is decreasing, $(C_n)$ is increasing. Apply part (a):
\[
\mathbb{P}\Big(\bigcup_{n\ge1}C_n\Big)=\lim_{n\to\infty}\mathbb{P}(C_n).
\]
Use complements:
\[
\Big(\bigcap_{n\ge1}B_n\Big)^c=\bigcup_{n\ge1}B_n^c=\bigcup_{n\ge1}C_n.
\]
Thus
\[
1-\mathbb{P}\Big(\bigcap_{n\ge1}B_n\Big)
=
\lim_{n\to\infty}\mathbb{P}(C_n)
=
\lim_{n\to\infty}(1-\mathbb{P}(B_n)),
\]
so $\mathbb{P}(\cap_{n\ge1}B_n)=\lim_{n\to\infty}\mathbb{P}(B_n)$.

\item \textbf{Example on $(0,1]$ with $A_n=(0,1/n]$.}

Here $(A_n)$ is \emph{decreasing} since $(0,1/(n+1)]\subset(0,1/n]$.
With uniform probability on $(0,1]$ (Lebesgue restricted and normalized),
\[
\mathbb{P}(A_n)=\lambda((0,1/n])=\frac{1}{n}.
\]
Moreover,
\[
\bigcap_{n\ge1}A_n=\bigcap_{n\ge1}(0,1/n]=\emptyset,
\]
since no $x>0$ can satisfy $x\le 1/n$ for all $n$.
Thus $\mathbb{P}(\cap_{n\ge1}A_n)=0$ and
\[
\lim_{n\to\infty}\mathbb{P}(A_n)=\lim_{n\to\infty}\frac1n=0,
\]
consistent with part (b).
\end{enumerate}

% -----------------------------
\item \textbf{Borel--Cantelli (I) and a concrete example — Solution}

Assume $\sum_{n=1}^\infty\mathbb{P}(A_n)<\infty$.

\begin{enumerate}[label=(\alph*)]
\item \textbf{Statement (BC I).}

\[
\mathbb{P}(\limsup_{n\to\infty}A_n)=0,
\quad\text{equivalently}\quad
\mathbb{P}(A_n\ \text{i.o.})=0.
\]

\item \textbf{Example with $\mathbb{P}(X_n=1)=1/n^2$.}

Let $A_n=\{X_n=1\}$. Then
\[
\mathbb{P}(A_n)=\frac{1}{n^2},
\qquad
\sum_{n=1}^\infty\mathbb{P}(A_n)=\sum_{n=1}^\infty\frac{1}{n^2}<\infty.
\]
By Borel--Cantelli (I),
\[
\mathbb{P}(A_n\ \text{i.o.})=\mathbb{P}(\limsup A_n)=0.
\]

\item \textbf{Interpretation.}

With probability $1$, only \emph{finitely many} indices $n$ satisfy $X_n=1$.
Equivalently, there exists a (random) $N(\omega)$ such that for all $n\ge N(\omega)$,
$X_n(\omega)=0$.
\end{enumerate}

% -----------------------------
\item \textbf{Borel--Cantelli (II) under independence — Solution}

Assume $(A_n)$ are independent and $\sum_{n=1}^\infty\mathbb{P}(A_n)=\infty$.

\begin{enumerate}[label=(\alph*)]
\item \textbf{Statement (BC II).}

If $(A_n)$ are independent and $\sum_n\mathbb{P}(A_n)=\infty$, then
\[
\mathbb{P}(\limsup_{n\to\infty}A_n)=1,
\quad\text{i.e.}\quad
\mathbb{P}(A_n\ \text{i.o.})=1.
\]

\item \textbf{Fair coin tosses, $A_n=\{X_n=H\}$.}

Here $\mathbb{P}(A_n)=1/2$ so $\sum_n\mathbb{P}(A_n)=\infty$ and the events are independent.
By BC II,
\[
\mathbb{P}(A_n\ \text{i.o.})=1,
\]
so with probability $1$ there are infinitely many heads.

\item \textbf{$A_n=\{X_n=H\ \text{and}\ X_{n+1}=H\}$ (two consecutive heads).}

These events are \emph{not} independent because $A_n$ and $A_{n+1}$ both involve $X_{n+1}$.
However, we can still conclude that $A_n$ occurs infinitely often with probability $1$:

Consider the \emph{disjoint} subsequence
\[
B_k := A_{2k-1}=\{X_{2k-1}=H,\ X_{2k}=H\}.
\]
The events $(B_k)$ are independent (they depend on disjoint pairs of tosses) and
\[
\mathbb{P}(B_k)=\frac14,
\qquad
\sum_{k=1}^\infty \mathbb{P}(B_k)=\infty.
\]
By BC II, $\mathbb{P}(B_k\ \text{i.o.})=1$.
But $B_k\subseteq A_n$ for infinitely many $n$ implies $A_n$ occurs infinitely often.
Hence $\mathbb{P}(A_n\ \text{i.o.})=1$.
\end{enumerate}

% -----------------------------
\item \textbf{Mixed distribution: censored exponential — Solution}

Let $Y\sim\mathrm{Exp}(1)$ and $X=\min\{1,Y\}$.

\begin{enumerate}[label=(\alph*)]
\item \textbf{Compute $\mathbb{P}(X=1)$.}

We have $X=1$ iff $Y\ge 1$. Thus
\[
\mathbb{P}(X=1)=\mathbb{P}(Y\ge1)=e^{-1}.
\]

\item \textbf{For $t\in[0,1)$, compute $F_X(t)$.}

If $t<1$, then $X\le t$ iff $Y\le t$ (since when $Y\ge1$, $X=1>t$). Hence
\[
F_X(t)=\mathbb{P}(X\le t)=\mathbb{P}(Y\le t)=1-e^{-t},
\qquad t\in[0,1).
\]
Also $F_X(t)=0$ for $t<0$.

\item \textbf{Density on $[0,1)$ and an atom at $1$.}

For $t\in(0,1)$,
\[
f_X(t)=F_X'(t)=e^{-t}.
\]
This is a density on $[0,1)$, with total mass
\[
\int_0^1 e^{-t}\,dt=1-e^{-1}.
\]
The remaining mass is at $1$:
\[
\mathbb{P}(X=1)=e^{-1},
\]
so $X$ has a mixed law (continuous part on $[0,1)$ plus an atom at $1$).

\item \textbf{Compute $\mathbb{E}[X]$.}

Using the hint and $\mathbb{P}(Y>t)=e^{-t}$,
\[
\mathbb{E}[X]=\mathbb{E}[\min\{1,Y\}]
=\int_0^1 \mathbb{P}(Y>t)\,dt
=\int_0^1 e^{-t}\,dt
=1-e^{-1}.
\]
(Equivalently: $\mathbb{E}[X]=\int_0^1 t e^{-t}\,dt + 1\cdot e^{-1}=1-e^{-1}$.)
\end{enumerate}

% -----------------------------
\item \textbf{Expectation may not exist (Cauchy) — Solution}

Let $f(x)=\frac{1}{\pi(1+x^2)}$.

\begin{enumerate}[label=(\alph*)]
\item \textbf{Show $\int_{\mathbb{R}} |x| f(x)\,dx = \infty$.}

By symmetry,
\[
\int_{\mathbb{R}} |x| f(x)\,dx
=2\int_0^\infty \frac{x}{\pi(1+x^2)}\,dx.
\]
Compute:
\[
\int_0^\infty \frac{x}{1+x^2}\,dx
=
\frac12\int_0^\infty \frac{2x}{1+x^2}\,dx
=
\frac12\Big[\ln(1+x^2)\Big]_{0}^{\infty}
=+\infty.
\]
Thus $\int_{\mathbb{R}} |x| f(x)\,dx=+\infty$.

\item \textbf{Deduce $\mathbb{E}[X]$ does not exist (Lebesgue).}

Lebesgue expectation $\mathbb{E}[X]$ exists only if
\[
\int_{\mathbb{R}} |x| f(x)\,dx<\infty.
\]
Since this integral is infinite, $X\notin L^1$ and $\mathbb{E}[X]$ does not exist.

\item \textbf{(Optional) Why symmetry does not guarantee existence.}

Although the density is symmetric, the positive and negative parts of $X$ have infinite expectation:
\[
\mathbb{E}[X^+]=\int_0^\infty x f(x)\,dx=+\infty,
\qquad
\mathbb{E}[X^-]=\int_{-\infty}^0 (-x) f(x)\,dx=+\infty.
\]
The expression ``$\infty-\infty$'' is undefined, so symmetry alone cannot define $\mathbb{E}[X]$.
\end{enumerate}

% -----------------------------
\item \textbf{Jensen: a useful bound — Solution}

Let $X\ge0$ be integrable.

\begin{enumerate}[label=(\alph*)]
\item \textbf{$\varphi(x)=x^2$.}

Since $x\mapsto x^2$ is convex,
\[
(\mathbb{E}[X])^2=\varphi(\mathbb{E}[X])\le \mathbb{E}[\varphi(X)]=\mathbb{E}[X^2].
\]

\item \textbf{$\varphi(x)=e^x$.}

Since $x\mapsto e^x$ is convex,
\[
e^{\mathbb{E}[X]}=\varphi(\mathbb{E}[X])\le \mathbb{E}[\varphi(X)]=\mathbb{E}[e^X].
\]

\item \textbf{Interpretation.}

For convex functions (like $\exp$), \emph{exponentiating then averaging} is larger than
\emph{averaging then exponentiating}:
\[
\mathbb{E}[e^X]\ \ge\ e^{\mathbb{E}[X]}.
\]
\end{enumerate}

% -----------------------------
\item \textbf{Joint density, marginals, and independence check — Solution}

Let
\[
f_{X,Y}(x,y)=6x^2y\,\mathbf{1}_{\{0\le x\le 1,\ 0\le y\le 1\}}.
\]

\begin{enumerate}[label=(\alph*)]
\item \textbf{Check normalization.}

\[
\int_0^1\int_0^1 6x^2y\,dy\,dx
=
\int_0^1 6x^2\Big[\frac{y^2}{2}\Big]_{0}^{1}\,dx
=
\int_0^1 3x^2\,dx
=
3\Big[\frac{x^3}{3}\Big]_{0}^{1}
=1.
\]

\item \textbf{Compute marginals.}

For $x\in[0,1]$,
\[
f_X(x)=\int_0^1 6x^2y\,dy
=6x^2\cdot \frac12
=3x^2.
\]
For $y\in[0,1]$,
\[
f_Y(y)=\int_0^1 6x^2y\,dx
=6y\Big[\frac{x^3}{3}\Big]_0^1
=2y.
\]

\item \textbf{Independence.}

For $(x,y)\in[0,1]^2$,
\[
f_X(x)f_Y(y)=(3x^2)(2y)=6x^2y=f_{X,Y}(x,y),
\]
so $X$ and $Y$ are independent.

\item \textbf{Compute $\mathbb{E}[X]$, $\mathbb{E}[Y]$, and $\mathrm{Var}(X)$.}

Using marginals:
\[
\mathbb{E}[X]=\int_0^1 x\cdot 3x^2\,dx = 3\int_0^1 x^3\,dx=3\cdot \frac14=\frac34.
\]
\[
\mathbb{E}[Y]=\int_0^1 y\cdot 2y\,dy = 2\int_0^1 y^2\,dy=2\cdot \frac13=\frac23.
\]
Also
\[
\mathbb{E}[X^2]=\int_0^1 x^2\cdot 3x^2\,dx=3\int_0^1 x^4\,dx=3\cdot \frac15=\frac35.
\]
Hence
\[
\mathrm{Var}(X)=\mathbb{E}[X^2]-(\mathbb{E}[X])^2
=\frac35-\Big(\frac34\Big)^2
=\frac35-\frac{9}{16}
=\frac{48-45}{80}
=\frac{3}{80}.
\]
\end{enumerate}

% -----------------------------
\item \textbf{Covariance and correlation basics — Solution}

Let $X,Y\in L^2$.

\begin{enumerate}[label=(\alph*)]
\item \textbf{Variance expansion.}

Write
\[
\mathrm{Var}(X+Y)=\mathbb{E}\big[(X+Y-\mathbb{E}[X+Y])^2\big]
=\mathbb{E}\big[((X-\mathbb{E}X)+(Y-\mathbb{E}Y))^2\big].
\]
Expand the square:
\[
\mathrm{Var}(X+Y)
=
\mathbb{E}[(X-\mathbb{E}X)^2]
+
\mathbb{E}[(Y-\mathbb{E}Y)^2]
+
2\mathbb{E}[(X-\mathbb{E}X)(Y-\mathbb{E}Y)].
\]
Thus
\[
\mathrm{Var}(X+Y)=\mathrm{Var}(X)+\mathrm{Var}(Y)+2\,\mathrm{Cov}(X,Y).
\]

\item \textbf{Independence implies variance additivity.}

If $X$ and $Y$ are independent and square-integrable, then
\[
\mathbb{E}[XY]=\mathbb{E}[X]\mathbb{E}[Y]
\quad\Rightarrow\quad
\mathrm{Cov}(X,Y)=0.
\]
Plugging into (a) gives
\[
\mathrm{Var}(X+Y)=\mathrm{Var}(X)+\mathrm{Var}(Y).
\]

\item \textbf{Counterexample: zero covariance but not independent.}

Let $X$ take values $\{-1,1\}$ with $\mathbb{P}(X=1)=\mathbb{P}(X=-1)=1/2$
(symmetric). Define $Y=X^2$.
Then $Y\equiv 1$ is constant, so $X$ and $Y$ are clearly not independent (indeed $Y$ is a function of $X$).
Compute covariance:
\[
\mathbb{E}[X]=0,\quad \mathbb{E}[Y]=1,\quad \mathbb{E}[XY]=\mathbb{E}[X\cdot 1]=\mathbb{E}[X]=0.
\]
Hence
\[
\mathrm{Cov}(X,Y)=\mathbb{E}[XY]-\mathbb{E}[X]\mathbb{E}[Y]=0-0\cdot 1=0,
\]
yet $X$ and $Y$ are not independent.
\end{enumerate}

% -----------------------------
\item \textbf{Two-asset minimum-variance portfolio}

Covariance matrix
\[
\Sigma=
\begin{pmatrix}
\sigma_1^2 & \rho\sigma_1\sigma_2\\
\rho\sigma_1\sigma_2 & \sigma_2^2
\end{pmatrix},
\qquad -1<\rho<1.
\]
Weights $\bm{w}=(w_1,w_2)$ with $w_1+w_2=1$.

\begin{enumerate}[label=(\alph*)]
\item Write $w_2=1-w_1$. Portfolio variance:
\begin{align*}
V(w_1)
&=
\bm{w}^\top\Sigma\bm{w}
=
\sigma_1^2 w_1^2 + \sigma_2^2(1-w_1)^2 + 2\rho\sigma_1\sigma_2 w_1(1-w_1)\\
&=
(\sigma_1^2+\sigma_2^2-2\rho\sigma_1\sigma_2)w_1^2
+(-2\sigma_2^2+2\rho\sigma_1\sigma_2)w_1+\sigma_2^2.
\end{align*}

\item Minimize the quadratic: $V'(w_1)=0$ gives
\[
2(\sigma_1^2+\sigma_2^2-2\rho\sigma_1\sigma_2)w_1+(-2\sigma_2^2+2\rho\sigma_1\sigma_2)=0,
\]
hence
\[
\boxed{\;
w_1^\star=\frac{\sigma_2^2-\rho\sigma_1\sigma_2}{\sigma_1^2+\sigma_2^2-2\rho\sigma_1\sigma_2},
\qquad
w_2^\star=1-w_1^\star=\frac{\sigma_1^2-\rho\sigma_1\sigma_2}{\sigma_1^2+\sigma_2^2-2\rho\sigma_1\sigma_2}.
\;}
\]

\item As $\rho\to -1$, the denominator increases and diversification becomes very powerful; in the limit $\rho=-1$ and $\sigma_1=\sigma_2$ one can achieve near-zero variance with $w_1=w_2=1/2$ (perfect hedging).
\end{enumerate}

\bigskip

% -----------------------------
\item \textbf{Finding a distribution via cdf (practice)}

$f_X(x)=2x\,\mathbf{1}_{[0,1]}(x)$, and $Y=\sqrt{X}$.

\begin{enumerate}[label=(\alph*)]
\item For $t<0$, $F_X(t)=0$. For $0\le t\le 1$,
\[
F_X(t)=\int_0^t 2x\,dx=t^2.
\]
For $t\ge 1$, $F_X(t)=1$.
So
\[
F_X(t)=
\begin{cases}
0,& t<0,\\
t^2,& 0\le t\le 1,\\
1,& t\ge 1.
\end{cases}
\]

\item For $t<0$, $F_Y(t)=0$. For $0\le t\le 1$,
\[
F_Y(t)=\mathbb{P}(Y\le t)=\mathbb{P}(\sqrt{X}\le t)=\mathbb{P}(X\le t^2)=F_X(t^2)=(t^2)^2=t^4.
\]
For $t\ge 1$, $F_Y(t)=1$.

\item Differentiate on $(0,1)$:
\[
f_Y(t)=F_Y'(t)=4t^3\,\mathbf{1}_{(0,1)}(t).
\]

\item Expectation:
\[
\mathbb{E}[Y]=\int_0^1 t\,f_Y(t)\,dt=\int_0^1 t\cdot 4t^3\,dt
=4\int_0^1 t^4\,dt=\frac{4}{5}.
\]
(Equivalently, $\mathbb{E}[\sqrt{X}]=\int_0^1 \sqrt{x}\,2x\,dx=2\int_0^1 x^{3/2}dx=4/5$.)
\end{enumerate}

\end{enumerate}

\end{document}